\section{Exact verification and ancillary material}\label{app:verification}

The five kinds of computer checks below use exact rational or
finite-field arithmetic; instructions, dependency versions and checksums
are included with the source package. The prime-index theorem does not
depend on the computer calculations.
\begin{enumerate}\setlength{\itemsep}{0pt}\setlength{\parskip}{0pt}
\item \texttt{verify\_manuscript.py} reconstructs the convolutions, the
first Laurent coefficients and the Weierstrass relations, and checks the
leading-term comparison in Proposition~\ref{prop:injective} and, by
exact substitution, the displayed formulas of orders two to four,
including the agreement of Table~\ref{tab:ks} with the tables of
\citet{KudryashovZargaryan1996,ConteMusette2009}, the multiplicities
\eqref{eq:ks-collisions}, the Kawahara coefficient restriction and its
literature normalization, the nodal correspondence and the Fisher
identity. The uniform proofs are in the text.
\item \texttt{p5\_count.sing} (\textsc{Singular}, Gr\"obner basis over
\(\QQ\)): at \(p=5\) the count with multiplicity of the full matching
equations, their quotient dimension, is \(126\). Adjoining the Jacobian
determinant gives the unit ideal: the determinant vanishes at no
solution, so all are simple.
\item At \(p=8\) the same two outputs, \(6435\) and the unit ideal, are
obtained over the finite field \(\FF_{32003}\). So every solution modulo
\(32003\) is simple and lifts uniquely by Hensel's lemma, as in Step~4
of the proof of Theorem~\ref{thm:prime}, and the lifts are all the
solutions because the count with multiplicity is \(6435\)
(Proposition~\ref{prop:general}(ii)). At \(p=11\) each of the \(352716\)
solutions is certified over the \(19\)-adic numbers, by Hensel's lemma
at the simple points modulo \(19\) and by exact Newton-polygon and
Newton--Kantorovich certificates at the singular ones. The archived
certificate and the record of its validation are in the repository
(\texttt{certificates/}).
\item \texttt{rigidity/} has the proof and exact certificate for
rigidity \((\mathrm R_d)\), \(d=p-1\le72\). For \(2d+1=\ell^e+g\),
\(g\in\{0,2,4,6\}\), the proof reduces the normalized continuous
convolution to a residual system of degree at most \(g\); the
prime-power case \(g=0\) is also in the companion note. The certificate
retains every residual degree, with integer ideal-membership identities
and certified prime factors; the replay multiplies out every identity
and checks all 72 arithmetic choices. It does not infer
characteristic-zero nonexistence from lists of residue-field points.
\item \texttt{symmetry\_checks/} checks Proposition~\ref{prop:reflection}
with symbolic coefficients for \(d\le5\); the labels of
Table~\ref{tab:ks} and of the \(21\) real pairs at \(p=4\), with
Proposition~\ref{cor:endpoints} and Theorem~\ref{thm:symmetric}(ii) for
them; the counts with multiplicity \(F_p\) for \(p\le8\), and
finite-field certificates, as at \(p=8\), that the \(F_p\) pairs are
distinct for every \(p\le14\); on a sample lattice, the counts with
multiplicity \(2F_p\) of Corollary~\ref{cor:elliptic-symmetric} for
\(p\le6\), and no elliptic pair of a purely dissipative equation for
\(p=3,5,7\) (\texttt{elliptic\_dissipative.py}); the six pairs
\eqref{eq:p5-symmetric} and the ten pairs \eqref{eq:p6-eliminant}, with
the sextic factor in full; and the scaling \eqref{eq:fixed-wave} by
substitution in the evolution equation.
\end{enumerate}

The Lean~4 formalization of Remark~\ref{rem:lean} is at
\url{https://github.com/migita/binomial-count-meromorphic-waves},
directory \texttt{lean/}, with Lean and Mathlib pinned to
version~4.33.1. It proves rigidity \((\mathrm R_d)\) for \(2d+1\) prime
(\texttt{rigidity\_prime\_index}) and the count and simplicity
statements of Theorem~\ref{thm:prime}, both for the matching equations
\eqref{eq:remainder} (\texttt{primeIndexCountAndSimplicity}) and for
the normalized pairs \((P,v)\) as solutions of \eqref{eq:main} in the
field of rational functions of \(e^{z}\)
(\texttt{primeIndexRationalExponentialCountAndSimplicity}), with the
correspondence between the two formulations for every \(p\ge2\). The
discrete convolution is defined there by the Bernoulli-polynomial power
sums; the remainder equations and their Jacobian are those of
Section~\ref{sec:counting}. \texttt{EnumerationCheck.lean} and
\texttt{Check.lean} verify that the proofs have exactly the stated
propositions as their types and depend only on the axioms
\texttt{propext}, \texttt{Classical.choice} and \texttt{Quot.sound}; the
project contains no \texttt{sorry}, added axiom or
\texttt{native\_decide}.
